% Derivazioni migrate dal Cap. 5 per il vincolo di lunghezza del corpo.
% Il capitolo mantiene il canale depolarizzante in forma completa e gli enunciati
% compatti degli altri tre.

I quattro canali quantistici che modellano il rumore sono mappe \ac{CPTP} in rappresentazione di Kraus, secondo l'Eq.~\ref{eq:kraus}. Il Capitolo~\ref{chap:rumore} ne riporta il canale depolarizzante per esteso, essendo quello su cui poggia la campagna sperimentale, e l'enunciato compatto degli altri tre; qui se ne danno le derivazioni complete.

\subsection{Canale Bit Flip}

Il canale \textbf{bit flip} modella un errore di tipo $X$ discreto: con probabilità $p$ il qubit viene invertito ($\ket{0}\leftrightarrow\ket{1}$), con probabilità $1-p$ rimane invariato. Gli operatori di Kraus sono $K_0 = \sqrt{1-p}\,I$ e $K_1 = \sqrt{p}\,X$, e il canale agisce come:
\begin{equation}
    \mathcal{E}_{BF}(\rho) = (1-p)\,\rho + p\,X\rho X
    \label{eq:bitflip}
\end{equation}
Sulla sfera di Bloch, le componenti $r_y$ e $r_z$ vengono contratte di un fattore $(1-2p)$ mentre la componente $r_x$ rimane invariata. Per $p=1/2$ il canale annulla le componenti $y$ e $z$, ma conserva $r_x$: equivale a una completa dephasing nella base degli autostati di $X$, non alla distruzione completa dello stato.

% ---
\subsection{Canale Phase Flip}

Il canale \textbf{phase flip} modella la perdita di informazione di fase senza cambiamento di popolazione: la sovrapposizione decade senza che l'energia venga dissipata. Con operatori di Kraus $K_0 = \sqrt{1-p}\,I$ e $K_1 = \sqrt{p}\,Z$, il canale è:
\begin{equation}
    \mathcal{E}_{PF}(\rho) = (1-p)\,\rho + p\,Z\rho Z
    \label{eq:phaseflip}
\end{equation}
L'effetto sulla matrice densità è immediato:
\begin{equation}
    \mathcal{E}_{PF}\!\begin{pmatrix}\rho_{00} & \rho_{01} \\ \rho_{10} & \rho_{11}\end{pmatrix}
    = \begin{pmatrix}\rho_{00} & (1-2p)\,\rho_{01} \\ (1-2p)\,\rho_{10} & \rho_{11}\end{pmatrix}
\end{equation}
Le popolazioni $\rho_{00}$ e $\rho_{11}$ rimangono invariate, mentre le \textbf{coerenze} $\rho_{01}$, $\rho_{10}$ vengono soppresse di un fattore $(1-2p)$. Questo è il canale che descrive il decadimento $T_2$: parametrizzando $p(t) = \frac{1}{2}(1-e^{-t/T_2})$, si ottiene $\rho_{01}(t) = \rho_{01}(0)\,e^{-t/T_2}$.

% ---
\subsection{Canale di Amplitude Damping}

Il canale di \textbf{amplitude damping} modella il rilassamento energetico $T_1$: il qubit decade spontaneamente dallo stato eccitato $\ket{1}$ allo stato fondamentale $\ket{0}$ con probabilità $\gamma = 1-e^{-t/T_1}$. Gli operatori di Kraus sono:
\begin{equation}
    K_0 = \begin{pmatrix}1 & 0 \\ 0 & \sqrt{1-\gamma}\end{pmatrix}, \qquad
    K_1 = \begin{pmatrix}0 & \sqrt{\gamma} \\ 0 & 0\end{pmatrix}
    \label{eq:kraus_ad}
\end{equation}
Il canale agisce sulla matrice densità come:
\begin{equation}
    \mathcal{E}_{AD}\!(\rho) = \begin{pmatrix}
        \rho_{00} + \gamma\,\rho_{11} & \sqrt{1-\gamma}\,\rho_{01} \\
        \sqrt{1-\gamma}\,\rho_{10}   & (1-\gamma)\,\rho_{11}
    \end{pmatrix}
    \label{eq:amp_damping_explicit}
\end{equation}
La popolazione $\rho_{11}$ decade verso $\rho_{00}$ con tasso $\gamma$, mentre le coerenze vengono soppresse di $\sqrt{1-\gamma}$. A differenza dei canali precedenti, questo canale è \textbf{non simmetrico}: sulla sfera di Bloch contrae la sfera verso il polo nord $\ket{0}$, non verso il centro:
\begin{equation}
    r_x \to \sqrt{1-\gamma}\,r_x, \qquad r_y \to \sqrt{1-\gamma}\,r_y, \qquad r_z \to (1-\gamma)\,r_z + \gamma
\end{equation}
Per $t \to \infty$ ($\gamma \to 1$), qualsiasi stato converge a $\ket{0}\bra{0}$, indipendentemente dallo stato iniziale.

\subsection{Confronto geometrico dei quattro canali}

La visualizzazione seguente, trasferita dal corpo, completa le derivazioni
mostrando che canali con probabilità d'errore simili possono produrre geometrie
molto diverse nello spazio degli stati.

Per il canale depolarizzante, l'identità
$X\rho X+Y\rho Y+Z\rho Z=2I-\rho$ (per $\operatorname{Tr}\rho=1$) conduce
dalla forma a miscela di Pauli alla forma compatta della
Eq.~\eqref{eq:depolarizing_compact}. Il vettore di Bloch è moltiplicato per
$1-4p/3$ e, per $p=3/4$, lo stato diventa massimamente misto qualunque fosse
l'ingresso. Nel parametro Aer $\lambda$, invece, la probabilità aggregata di
una Pauli non identità è $3\lambda/4$ su un qubit e $15\lambda/16$ su due.
I preset uniformi compongono depolarizzazione e thermal relaxation; M8 usa il
solo proxy Pauli per gate, mentre M11 converte l'infedeltà media della snapshot
senza aggiungere nuovamente $T_1/T_2$, evitando il doppio conteggio.

\begin{figure}[H]
\centering
\includegraphics[width=0.7\textwidth]{14_bloch_channels.pdf}
\caption[Effetto dei canali di rumore sulla sfera di Bloch]{Effetto dei quattro
canali nel piano $xz$. La curva tratteggiata è la sfera ideale: bit flip e
phase flip contraggono direzioni selettive, il depolarizzante contrae
isotropicamente e l'amplitude damping sposta lo stato verso $\ket{0}$.}
\label{fig:bloch_channels}
\end{figure}

\begin{table}[H]
\centering
\small
\begin{tabular}{@{}llll@{}}
\toprule
\textbf{Canale} & \textbf{Effetto sulle coerenze} & \textbf{Popolazioni} & \textbf{Modella} \\
\midrule
Bit flip & contrae $r_y,r_z$ & scambiate & errore $X$ \\
Phase flip & sopprime $\rho_{01}$ & invariate & dephasing $T_2$ \\
Depolarizzante & contrazione isotropa & verso $I/2$ & gate generico \\
Amplitude damping & sopprime $\rho_{01}$ & verso $\ket{0}$ & rilassamento $T_1$ \\
\bottomrule
\end{tabular}
\caption{Confronto fra i quattro canali fondamentali~\cite{nielsenChuang2000}.}
\label{tab:canali_rumore}
\end{table}

% La sezione seguente genera intenzionalmente A.5 pur essendo collocata in
% coda al file dei canali: evita di modificare la struttura di main.tex.
\section{Strategie anti-rumore e richiami estesi di QEC}
\label{app:strategie_estese}

Questa sezione raccoglie il materiale trasferito dalla parte del
Capitolo~\ref{chap:quadro} dedicata alle strategie. Nel corpo restano la tabella
comparativa, il ciclo QEC essenziale e le scelte che guidano gli esperimenti;
qui sono conservati rassegna, esempi e confronti estesi.

\subsection{I quattro livelli di intervento}

\begin{enumerate}
    \item \textbf{Soppressione del rumore.} Agisce prima o durante il calcolo
    riducendo fisicamente l'accoppiamento con l'ambiente. Comprende
    \emph{dynamical decoupling}, ottimizzazione della durata dei gate e
    isolamento criogenico ed elettromagnetico. Riduce il rumore ma è vincolata
    all'hardware disponibile.

    \item \textbf{Mitigazione degli errori.} Opera a posteriori senza
    correggere i qubit. \ac{ZNE}, \emph{probabilistic error cancellation} e
    gate twirling combinano più esecuzioni o randomizzazioni per stimare
    l'output ideale. Il limite è l'overhead di campionamento, la cui legge
    dipende dalla tecnica e può diventare esponenziale nell'errore totale,
    soprattutto per la PEC.

    \item \textbf{Tolleranza ai guasti.} Codifica ogni qubit logico in molti
    qubit fisici. Sotto le ipotesi del threshold theorem e al di sotto della
    soglia~\cite{aharonov1997fault}, l'errore logico può essere ridotto
    aumentando distanza e overhead. Le stime per algoritmi di grande scala
    possono richiedere centinaia o migliaia di qubit fisici per qubit logico,
    a seconda di codice, errore fisico e obiettivo.

    \item \textbf{Metodi appresi.} Modelli classici possono classificare
    output, stimare qualità o assistere un decoder. Non sono automaticamente
    \ac{QML}, né sostituiscono per definizione mitigazione o QEC.
\end{enumerate}

Nel perimetro della tesi la soppressione non è controllabile, la mitigazione
richiede esecuzioni aggiuntive e la fault tolerance completa eccede le risorse
disponibili. Un modello a valle non modifica il circuito, ma il suo vantaggio è
un'ipotesi da confrontare con regole semplici. Per questo TOP-1, TOP-4 e M2
usano gli stessi istogrammi, mentre la \ac{ZNE} è un controllo secondario.

\subsection{Usi possibili dei metodi appresi}

Il termine appropriato per gli esperimenti è \textbf{machine learning classico
applicato a dati quantistici}: Random Forest, \ac{SVM} e \ac{MLP} sono
addestrati ed eseguiti classicamente. Il \ac{QML} in senso stretto impiega
circuiti o modelli quantistici nell'apprendimento e non è implementato. Tre usi
distinti chiariscono il perimetro:
\begin{itemize}
    \item nella \textbf{caratterizzazione del rumore}, una regressione può
    collegare $T_1,T_2$, profondità e fedeltà per predire la qualità;
    \item nell'\textbf{ottimizzazione circuitale}, un agente può adattare
    decomposizione e routing al profilo d'errore del dispositivo;
    \item nella \textbf{decodifica QEC}, un modello usa le sindromi per
    proporre una correzione o correggere residualmente una decisione analitica.
\end{itemize}
Il terzo caso colloca il machine learning \emph{dentro una fase} della QEC,
senza ridefinire la disciplina.

\subsection{Principio, codici e decoder QEC}
\label{app:qec_richiami_teoria}

Il no-cloning impedisce di proteggere un qubit duplicandolo, ma non di
\emph{distribuire} l'informazione logica su più qubit fisici. Misurando
operatori collettivi si ottiene una sindrome dell'errore senza rivelare, e
quindi distruggere, lo stato logico. Il paradigma comprende:
\begin{enumerate}
    \item codifica di $\ket{\psi_L}$ su $n$ qubit fisici;
    \item misura di stabilizzatori che commutano con i codeword e
    anticommutano con determinate classi d'errore;
    \item decodifica della sindrome e scelta dell'operazione di recupero.
\end{enumerate}

Il \textbf{codice di Shor}~\cite{shorcorrection1995} codifica un qubit in nove
e corregge un errore singolo arbitrario:
\begin{equation}
 \ket{0_L}=\frac{(\ket{000}+\ket{111})^{\otimes3}}{2\sqrt2},\qquad
 \ket{1_L}=\frac{(\ket{000}-\ket{111})^{\otimes3}}{2\sqrt2}.
\end{equation}
La struttura combina protezione dalla fase mediante interferenza e dal bit flip
mediante ripetizione. Il \textbf{codice di Steane}~\cite{steane1996} usa sette
qubit e la struttura del codice classico di Hamming $[7,4,3]$; è un codice
\ac{CSS}~\cite{calderbankShor1996} e corregge anch'esso un errore singolo. Il
\textbf{surface code}~\cite{fowler2012} dispone invece i qubit su una griglia
2D con interazioni locali: sotto una soglia dipendente da circuito, rumore e
decoder, aumentare la distanza $d$ sopprime l'errore logico; sopra soglia lo
amplifica.

Il punto di contatto con il machine learning è la decodifica
\cite{torlaiMelko2017,alphaqubit2024}. Un modello può sostituire il decoder o
intervenire sui suoi residui, ma deve essere confrontato a parità di sindromi e
calendario. Nel protocollo studiato la sostituzione del matching non conviene;
l'affiancamento recupera invece una parte degli errori correlati assenti dal
grafo analitico. Questo è un risultato circoscritto al modello simulato, non una
superiorità generale dei decoder appresi.

\subsection{Verifica del filtro appreso nella prima fase}

Senza ML, ogni iterazione produce un istogramma: TOP-1 prova la moda e TOP-4 i
primi quattro esiti. Le iterazioni non vanno confuse con gli shot che formano
ciascun istogramma; nelle campagne la media è fra 1 e 1.5 iterazioni. Con ML,
il classificatore riceve le feature dell'istogramma già campionato e decide se
attivare TOP-4: non classifica il singolo shot e non elimina il costo di
costruzione della distribuzione.

L'ipotesi preregistrata era una riduzione rispetto a TOP-1. Essa si osserva
(nello scenario moderato da $\bar M_1=1.37$ a $\bar M=1.00$), ma l'ablazione
mostra che deriva dalla ricerca multi-candidato: il classificatore costa più di
TOP-4 in entrambi gli scenari. Il risultato negativo delimita quel filtro e
quelle etichette; non è una conclusione generale su ML, mitigazione o decoder.

\subsection{AQFT e scambio fra qubit e porte}
\label{app:aqft_estesa}

La \ac{QFT} impiega rotazioni controllate con angolo esponenzialmente più
piccolo all'aumentare della distanza fra qubit. L'\ac{AQFT} scarta quelle sotto
una soglia: perde precisione algoritmica, ma riduce il conteggio di porte. Per
$N=21$, troncare le trasformate interne agli addizionatori porta le porte a due
qubit da 49\,661 a 23\,178; nelle tre dimensioni del registro di phase
estimation la variante troncata supera quella completa in tutti i nove
confronti. Ridurre la profondità è una strategia anti-rumore perché accorcia il
cammino critico rispetto a $T_1$ e $T_2$, indipendentemente dal numero di qubit
disponibili.

Chevignard, Fouque e Schrottenloher~\cite{chevignard2025} perseguono uno scambio
diverso: l'aritmetica modulare approssimata in un sistema di residui riduce i
qubit logici da circa $2n$ a $n/2+o(n)$; per RSA-2048 stima 1730 qubit contro i
6144 del riferimento. Gli autori riportano però $2^{35.5}$ porte Toffoli e in
media quaranta esecuzioni, $2^{40.9}$ complessive, contro $2^{31.3}$ per una
singola esecuzione di riferimento. Un lavoro successivo~\cite{gidney2025}
riduce il costo e arriva a 1399 qubit logici. Nel regime fault-tolerant domina
il costo dei qubit; nel regime NISQ simulato qui domina la profondità rispetto
alla coerenza. Le due strategie rispondono quindi a vincoli diversi.

\subsection{Transizione fra le due fasi sperimentali}

Il primo esito delimita il post-processing: l'istogramma contiene candidati
utili, sfruttati da TOP-4, ma il classificatore su quelle etichette non aggiunge
informazione. Il nucleo successivo si sposta durante il calcolo: codice a
ripetizione, Steane, surface code e confronto fra decodifica analitica e
appresa. L'analisi fenomenologica di Shor resta separata dai tassi QEC per
ciclo, mentre l'\ac{AQFT} valuta una leva circuitale complementare.

\subsection{Dettagli tecnici del contratto sperimentale}
\label{app:contratto_metodo_esteso}

Questo blocco conserva i dettagli operativi alleggeriti nella sezione
metodologica. Il circuito è compilato una sola volta nella base
\texttt{rz/sx/x/cx}, con \texttt{optimization\_level=2} e seed del transpiler
fissato. La configurazione di riferimento contiene 224 porte \texttt{cx} e ha
profondità 412; hash e conteggi completi della netlist sono registrati nel
manifest. Per ogni coppia (replica, iterazione), Aer genera un unico istogramma
che viene consegnato senza ricampionamento a M1, TOP-4 e M2.

I pareggi di frequenza sono risolti mediante una priorità SHA-256 dipendente
dal seed e dalla bitstring. In questo modo l'ordine dei candidati non dipende
né dall'ordine interno del dizionario né dal valore numerico dell'esito. Il
dataset conserva sia l'etichetta TOP-1 sia l'audit TOP-16, mentre la pipeline
usa esclusivamente il modello TOP-1 compatibile con il manifest. Ciascuno
scenario contiene 30 repliche e un budget di 50 iterazioni; il fallimento al
budget assume la sentinella 51. I test inferenziali agiscono sempre sui vettori
appaiati completi, non sulle sole medie aggregate.
